%
% $Id: slides.tex 4228 2010-12-02 17:55:12Z pcoca $
%
%
% Compilar a .pdf con LaTeX (pdflatex)
% Es necesario instalar Beamer (paquete latex-beamer en Debian)
%

%
% Gráficos:
% Los gráficos pueden suministrarse en PNG, JPG, TIF, PDF, MPS
% Los EPS deben convertirse a PDF (usar epstopdf)
%

\documentclass{beamer}
\usetheme{Warsaw}
\usebackgroundtemplate{\includegraphics[width=\paperwidth]{format/libresoft-bg.png}}
\usepackage[spanish]{babel}
\usepackage[utf8]{inputenc}
\usepackage{graphics}
\usepackage{amssymb} % Simbolos matematicos
\usepackage{url}

%\definecolor{libresoftgreen}{RGB}{162,190,43}
%\definecolor{libresoftblue}{RGB}{0,98,143}

%\setbeamercolor{titlelike}{bg=libresoftgreen}

%% Metadatos del PDF.
\hypersetup{
  pdftitle={FLOSS Testing and Quality Assurance},
  pdfauthor={F. Ortega, D. Izquierdo, P. Coca, P. García}
  pdfcreator={GSyC/Libresoft},
  pdfproducer=PDFLaTeX,
  pdfsubject={nn},
}
%%

\AtBeginSection[]
{
  \begin{frame}<presentation>
    \frametitle{Index}
    \tableofcontents[current]
  \end{frame}
}


\begin{document}

\title{FLOSS Testing and Quality Assurance}
\subtitle{Mozilla Thunderbird}
\institute{\\pcoca@libresoft.es\\
GSyC/Libresoft}
\author{Felipe Ortega, Daniel Izquierdo, Pedro Coca, Pedro García}
\date{\today}

\frame{
\maketitle
\begin{center}
\includegraphics[width=6cm]{format/gsyc-urjc}
\end{center}
}


% Si el titulo o el autor se quieren acortar para los pies de página
% se pueden redefinir aquí:
%\title{Titulo corto}
%\author{Autores abreviado}


%% LICENCIA DE REDISTRIBUCION DE LAS TRANSPAS
\frame{
~
\vspace{4cm}

\begin{flushright}
{\tiny
(cc) 2010 Felipe Ortega, Daniel Izquierdo, Pedro Coca, Pedro García. \\
Some rights reserved. This document is distributed under the Creative \\
            Commons Attribution-ShareAlike 3.0 licence, available in \\
            http://creativecommons.org/licenses/by-sa/3.0/

%  Este documento (o uno muy similar) está disponible en \\
%  \url{http://gsyc.escet.urjc.es/~jjamor/}
}
\end{flushright}
}
%%

%%%%%%
%Transpas separadas por \begin{frame}
%%%%%%%%%%%%%%%%%%%%%%%%\end{frame}

\section{Introduction}

\begin{frame}
\frametitle{What is QA}
\begin{itemize}
\item Tracking and monitoring the software engineering processes
    \begin{itemize}
       \item software design 
       \item source code control 
       \item change management 
       \item ...
       \item configuration management 
       \item testing 
       \item release management
    \end{itemize}
\item Goal: Ensure "Quality"
\item As we have seen before, some standards, methodologies and paradigms aim partially or totally this goal
\end{itemize}
\end{frame}

%%%%%%%%%%%%%%%%%%%%%%%%%%%%%%%%%%%%%%%%%%%%%%%%%%%%%%%%%%%%%%

\section{Introduction to Mozilla}

%%%%%%%%%%%%%%%%%%%%%%%%%%%%%%%%%%%%%%%%%%%%%%%%%%%%%%%%%%%%%%

\begin{frame}
\frametitle{What is Mozilla?}
 \begin{itemize}
 \item What do you know about Mozilla project?
 \end{itemize}
\end{frame}

%%%%%%%%%%%%%%%%%%%%%%%%%%%%%%%%%%%%%%%%%%%%%%%%%%%%%%%%%%%%%%


\begin{frame}
\frametitle{What is Mozilla?}
 \begin{itemize}
 \item The Mozilla Foundation is a non-profit organization that exists to support and provide leadership for the Open Source Mozilla project
 \item The Mozilla Foundation is formed by two subsidiaries:
     \begin{itemize}
     \item Mozilla Corporation. Mozilla Firefox web browser
     \item Mozilla Messaging. Mozilla Thunderbird email client
     \end{itemize}

 \end{itemize}
\end{frame}

%%%%%%%%%%%%%%%%%%%%%%%%%%%%%%%%%%%%%%%%%%%%%%%%%%%%%%%%%%%%%%

\begin{frame}
\frametitle{Some history about Mozilla}
 \begin{itemize}
 \item Netscape Communications Corporation
 \item Internal codename for the Netscape Navigator (Web Browser)
    \begin{itemize}
    \item Mosaic Killer. Project pet (Final logo defined in 1994)
    \end{itemize}
 \item Mozilla Application Suite (Now Seamonkey)
 \item Mozilla Organisation (1998)
 \item Mozilla Foundation (2003)
 \item Mozilla Corporation (2005) For-profit taxable subsidiary of the foundation
 \item Mozilla Foundation holds the Mozilla trademarks (2006)
 \item Mozilla Messaging (2008) For-profit taxable subsidiary of the foundation
 \end{itemize}
\end{frame}

%%%%%%%%%%%%%%%%%%%%%%%%%%%%%%%%%%%%%%%%%%%%%%%%%%%%%%%%%%%%%%

\begin{frame}
\frametitle{Some facts about Mozilla Foundation}
 \begin{itemize}
 \item Based on Mountain View (California, USA)
 \item Defined as "non-profit organization dedicated to preserving choice and promoting innovation on the Internet"
 \item 91,3 USD million of Revenue (2009) 
 \item 4 employees
 \item Mozilla Foundation has 3 subsidiaries (Corporation, Messaging, Online)
 \end{itemize}
\end{frame}

%%%%%%%%%%%%%%%%%%%%%%%%%%%%%%%%%%%%%%%%%%%%%%%%%%%%%%%%%%%%%%

\begin{frame}
\frametitle{Some facts about Mozilla Corporation}
 \begin{itemize}
 \item Based on Mountain View (California, USA). With offices in Toronto, Paris, Tokyo, etc.
 \item 91,3 USD million of Revenue (2009) 
 \item More than 250 employees
 \item 1 subsidiary (Mozilla Online, China)
 \end{itemize}
\end{frame}

%%%%%%%%%%%%%%%%%%%%%%%%%%%%%%%%%%%%%%%%%%%%%%%%%%%%%%%%%%%%%%

\begin{frame}
\frametitle{Some facts about Mozilla Messaging}
 \begin{itemize}
 \item Based on Vancouver, Canada
 \item AKA MoMo
 \item 10 employees
 \end{itemize}
\end{frame}

%%%%%%%%%%%%%%%%%%%%%%%%%%%%%%%%%%%%%%%%%%%%%%%%%%%%%%%%%%%%%%

\section{Mozilla QA}

%%%%%%%%%%%%%%%%%%%%%%%%%%%%%%%%%%%%%%%%%%%%%%%%%%%%%%%%%%%%%%

\begin{frame}
\frametitle{What is QA for Mozilla?}
 \begin{itemize}

 \item Quality Assurance is the method of systematically evaluating and monitoring various aspects of a project. Mozilla does this by using various tools such as Bugzilla, Crash Reporter, Manual Software Testing, Automated Software Testing and many more!

 \item In general, Software Testing at Mozilla is treated as an objective inquiry conducted by QA Community to provide the project with information about the quality of Mozilla software, feature-set and usability from the perspective of the user.

 \end{itemize}
\end{frame}

%%%%%%%%%%%%%%%%%%%%%%%%%%%%%%%%%%%%%%%%%%%%%%%%%%%%%%%%%%%%%%

\begin{frame}
\frametitle{Mozilla QA Teams}
 \begin{itemize}
 \item Desktop Firefox QA
 \item Browser Technologies QA
 \item Web QA
 \item QA Test automation
 \item QA Services
 \end{itemize}
\end{frame}

%%%%%%%%%%%%%%%%%%%%%%%%%%%%%%%%%%%%%%%%%%%%%%%%%%%%%%%%%%%%%%

\begin{frame}
\frametitle{Mozilla QA structure}
\begin{center}
 \includegraphics[height=5cm,width=11.5cm]{figs/MozillaQAOrgPic.jpg}
\begin{figure}
\end{figure}
\end{center}
\end{frame}

%%%%%%%%%%%%%%%%%%%%%%%%%%%%%%%%%%%%%%%%%%%%%%%%%%%%%%%%%%%%%%

\begin{frame}
\frametitle{Mozilla QA Teams}
 \begin{itemize}
 \item Desktop Firefox QA
      \begin{itemize}
      \item Performs testing and release assessment of the desktop FF browser
      \item Focus for upcoming major releases and maintenance releases
      \item Anyone can participate!
      \end{itemize}
 \end{itemize}
\end{frame}

%%%%%%%%%%%%%%%%%%%%%%%%%%%%%%%%%%%%%%%%%%%%%%%%%%%%%%%%%%%%%%

\begin{frame}
\frametitle{Mozilla QA Teams}
 \begin{itemize}
 \item Browser Technologies QA
      \begin{itemize}
      \item Performs testing and release assessment of other FF related technologies and services.
      \item Mobile and Web services are becoming more popular, therefore more prominent in the Mozilla project.
      \item Test Pilot, Firefox Syn, Jetpack, Mobile Firefox...
      \item Anyone can participate!
      \end{itemize}
 \end{itemize}
\end{frame}

%%%%%%%%%%%%%%%%%%%%%%%%%%%%%%%%%%%%%%%%%%%%%%%%%%%%%%%%%%%%%%

\begin{frame}
\frametitle{Mozilla QA Teams}
 \begin{itemize}
 \item Web QA
      \begin{itemize}
      \item Performs testing and release assessment of the Mozilla web sites
      \item Internal and web applications
      \item Several tools, such as Feed/HTML5/W3C/etc validators, XSS and SQL injection checking, performance checks, etc
      \item Anyone can participate!
      \end{itemize}
 \end{itemize}
\end{frame}

%%%%%%%%%%%%%%%%%%%%%%%%%%%%%%%%%%%%%%%%%%%%%%%%%%%%%%%%%%%%%%

\begin{frame}
\frametitle{Mozilla QA Teams}
 \begin{itemize}
 \item QA Test Automation
      \begin{itemize}
      \item Creates automated test and tools to support automation testing
      \item Uses Mozmill test tool
      \item Automate existing manual test (added to Litmus)
      \item Aims to lower the active testing time for QA engineers and contributors
      \item Anyone can participate!
      \end{itemize}
 \end{itemize}
\end{frame}

%%%%%%%%%%%%%%%%%%%%%%%%%%%%%%%%%%%%%%%%%%%%%%%%%%%%%%%%%%%%%%

\begin{frame}
\frametitle{Mozilla QA Teams}
 \begin{itemize}
 \item QA Services
      \begin{itemize}
      \item Provide services to support the QA community
      \item Primarily end user focused
      \item CrashKill
         \begin{itemize}
         \item Triage Top Crashes
         \item Identify top tools issues
         \item Identify automation priorities
         \end{itemize}
      \item Feedback review
      \item Accessibility technologies
      \item Bug Days, Test Days, Community Plans, etc
      \item Anyone can participate!
      \end{itemize}
 \end{itemize}

\end{frame}

%%%%%%%%%%%%%%%%%%%%%%%%%%%%%%%%%%%%%%%%%%%%%%%%%%%%%%%%%%%%%%

\begin{frame}
\frametitle{QA Teams: Anyone can participate}
 \begin{itemize}
    \item File bugs
    \item Triage bugs (confirm existing bugs and assign them to the right buckets)
    \item Help test new features
    \item Write test cases
    \item Plan new features testing
    \item Help others who want to get involved. 
 \end{itemize}
\end{frame}

%%%%%%%%%%%%%%%%%%%%%%%%%%%%%%%%%%%%%%%%%%%%%%%%%%%%%%%%%%%%%%

\begin{frame}
\frametitle{Mozilla QA Community}
\begin{center}
 \includegraphics[height=6cm]{figs/MozillaTestDay.png}
\begin{figure}
\end{figure}
\end{center}
\end{frame}


%%%%%%%%%%%%%%%%%%%%%%%%%%%%%%%%%%%%%%%%%%%%%%%%%%%%%%%%%%%%%%

\section{Introduction to Thunderbird}

%%%%%%%%%%%%%%%%%%%%%%%%%%%%%%%%%%%%%%%%%%%%%%%%%%%%%%%%%%%%%%

\begin{frame}
\frametitle{What is Thunderbird?}
 \begin{itemize}
 \item Thunderbird is a Free/Libre/Open Source Cross-Platform email and news client
 \item Runs on Windows, Linux, MacOS X, OS/2, OpenSolaris, etc
 \end{itemize}
\end{frame}


%%%%%%%%%%%%%%%%%%%%%%%%%%%%%%%%%%%%%%%%%%%%%%%%%%%%%%%%%%%%%%

\begin{frame}
 \frametitle{Thunderbird Features}
 \begin{itemize}
 \item Message management
 \item Tabs
 \item Mail Account Setup Wizard
 \item One-click Address Book
 \item Attachment Reminder
 \item Activity Manager
 \item Offline operation
 \item Message Search
 \item Add-ons (Calendar, etc)
 \end{itemize}
\end{frame}

%%%%%%%%%%%%%%%%%%%%%%%%%%%%%%%%%%%%%%%%%%%%%%%%%%%%%%%%%%%%%%

\begin{frame}
 \frametitle{Thunderbird Licenses}
 \begin{itemize}
 \item Mozilla is released under several FLOSS Licences
 \item Triple licensing:
    \begin{itemize}
     \item LGPL
     \item GPLv2
     \item Mozilla Public License 1.1
    \end{itemize}

 \end{itemize}
\end{frame}

%%%%%%%%%%%%%%%%%%%%%%%%%%%%%%%%%%%%%%%%%%%%%%%%%%%%%%%%%%%%%%

\begin{frame}
 \frametitle{Versions}
 \begin{itemize}
 \item Initial release
    \begin{itemize}
     \item 0.1. July 2003
    \end{itemize}

 \item Current stable release (December 2010)
    \begin{itemize}
     \item 3.1.6. October 2010
    \end{itemize}

 \item Latest alpha release (December 2010)
    \begin{itemize}
     \item 3.3.a1 November 2010
    \end{itemize}

 \end{itemize}
\end{frame}

%%%%%%%%%%%%%%%%%%%%%%%%%%%%%%%%%%%%%%%%%%%%%%%%%%%%%%%%%%%%%%

\begin{frame}
 \frametitle{Versions}
 \begin{itemize}
 \item First and Second digit means major releases
 \item Third digit means security and stability releases
 \item Initial release
    \begin{itemize}
     \item 0.1. July 2003
    \end{itemize}

 \item First mayor release
    \begin{itemize}
     \item 1.0. December 2004
    \end{itemize}

 \item Second mayor release
    \begin{itemize}
     \item 2.0. April 2007
    \end{itemize}

 \item Third mayor release
    \begin{itemize}
     \item 3.0. December 2009
    \end{itemize}

 \end{itemize}
\end{frame}

%%%%%%%%%%%%%%%%%%%%%%%%%%%%%%%%%%%%%%%%%%%%%%%%%%%%%%%%%%%%%%

\begin{frame}
\frametitle{Software Metrics}
 \begin{itemize}
 \item About 1.000.000 lines of code
 \item Programming Languages
    \begin{itemize}
     \item C++
     \item Javascript
     \item XML
     \item Java
     \item C
     \item Others
    \end{itemize}
 \item Almost 600 contributors
 \item 6 Million daily users
 \end{itemize}
\end{frame}

%%%%%%%%%%%%%%%%%%%%%%%%%%%%%%%%%%%%%%%%%%%%%%%%%%%%%%%%%%%%%%

\begin{frame}
 \frametitle{Mozilla Thunderbird QA}
 \begin{itemize}
 \item Mozilla Thunderbird QA uses tools such as:
     \begin{itemize}
     \item Bugzilla
     \item Crash reporter
     \item Manual software testing
     \item Automated software testing 
     \end{itemize}
 \end{itemize}
\end{frame}

%%%%%%%%%%%%%%%%%%%%%%%%%%%%%%%%%%%%%%%%%%%%%%%%%%%%%%%%%%%%%%

\section{Bug Reporting}

%%%%%%%%%%%%%%%%%%%%%%%%%%%%%%%%%%%%%%%%%%%%%%%%%%%%%%%%%%%%%%

\begin{frame}
 \frametitle{Bug Reporting}
 \begin{itemize}

    \item Once reported, Developers study in detail the bug information
    \item The information in a bug report could hints at the cause
    \item The more info, the better!
    \item A bug report without useful information could be just a waste of time: 
    \item Information provided in bug reports affects
         \begin{itemize}
         \item Resolution effort
         \item Resolution time
         \end{itemize}

 \end{itemize}
\end{frame}

%%%%%%%%%%%%%%%%%%%%%%%%%%%%%%%%%%%%%%%%%%%%%%%%%%%%%%%%%%%%%%

\begin{frame}
 \frametitle{Bug Writing Guidelines }
 \begin{itemize}
    \item Be precise
    \item Be clear: explain it so others can reproduce the bug
    \item One bug per report
    \item No bug is too trivial to report: small bugs may hide big bugs
    \item Clearly separate fact from speculation
 \end{itemize}
\end{frame}

%%%%%%%%%%%%%%%%%%%%%%%%%%%%%%%%%%%%%%%%%%%%%%%%%%%%%%%%%%%%%%

\begin{frame}
 \frametitle{Bad Bug Report examples}
 \begin{itemize}
     \item Browser should work with my web site
     \item I just clicked and it crashes
     \item System is really slow
     \item ... happens sometimes
 \end{itemize}
\end{frame}

%%%%%%%%%%%%%%%%%%%%%%%%%%%%%%%%%%%%%%%%%%%%%%%%%%%%%%%%%%%%%%

\begin{frame}
 \frametitle{Bad Bug Report examples}
\begin{center}
 \includegraphics[height=6cm]{figs/BadBugEmail.png}
\begin{figure}
\end{figure}
\end{center}
\end{frame}

%%%%%%%%%%%%%%%%%%%%%%%%%%%%%%%%%%%%%%%%%%%%%%%%%%%%%%%%%%%%%%

\begin{frame}
 \frametitle{What Makes a Good Bug Report?}
 \begin{itemize}
    \item Survey conducted. 150 developers and reporters
        \begin{itemize}
        \item Apache
        \item Eclipse
        \item Mozilla
        \end{itemize}
 \end{itemize}
\end{frame}

%%%%%%%%%%%%%%%%%%%%%%%%%%%%%%%%%%%%%%%%%%%%%%%%%%%%%%%%%%%%%%

\begin{frame}
 \frametitle{Bug reporting empirical studies: Bug reports contents}
 \begin{itemize}
    \item Questions for the Developers:
      \begin{itemize}
      \item 1) Items previously used when fixing bugs
      \item 2) Top 3 helpful items 
      \item 16 Options
      \end{itemize}
 \end{itemize}
\end{frame}

%%%%%%%%%%%%%%%%%%%%%%%%%%%%%%%%%%%%%%%%%%%%%%%%%%%%%%%%%%%%%%

\begin{frame}
 \frametitle{Bug reporting empirical studies: Bug reports contents}
 \begin{itemize}
    \item Developers responses about more valuable information: 
        \begin{itemize}
        \item How to reproduce the bug (steps)
        \item Stack traces 
        \item Test cases
        \item Screenshots (Essentials for GUI testing!)
        \item Code examples
        \end{itemize}
 \end{itemize}
\end{frame}

%%%%%%%%%%%%%%%%%%%%%%%%%%%%%%%%%%%%%%%%%%%%%%%%%%%%%%%%%%%%%%

\begin{frame}
 \frametitle{Bug reporting empirical studies: Bug reports contents}
 \begin{itemize}
    \item Questions for the Developers:
      \begin{itemize}
      \item 1) Issues found when fixing bugs
      \item 2) Top 3 delaying issues when fixing bugs
      \item 21 Options
      \end{itemize}
 \end{itemize}
\end{frame}

%%%%%%%%%%%%%%%%%%%%%%%%%%%%%%%%%%%%%%%%%%%%%%%%%%%%%%%%%%%%%%

\begin{frame}
 \frametitle{Bug reporting empirical studies: Bug reports contents}
 \begin{itemize}
    \item Developers responses about problem's severeness
        \begin{itemize}
        \item How to reproduce the bug (steps)
        \item Incomplete information
        \item Observed behavior
        \item Test cases
        \end{itemize}
 \end{itemize}
\end{frame}

%%%%%%%%%%%%%%%%%%%%%%%%%%%%%%%%%%%%%%%%%%%%%%%%%%%%%%%%%%%%%

\begin{frame}
 \frametitle{Bug reporting empirical studies: Bug reports contents}
 \begin{itemize}
    \item Questions for the Reporters:
      \begin{itemize}
      \item 1) Items reported in the past
      \item 2) Top 3 more difficult items
      \item 2) Top 3 items considered most relevant for developers
      \item Same 16 items offered to the previous group (developers)
      \end{itemize}
 \end{itemize}
\end{frame}

%%%%%%%%%%%%%%%%%%%%%%%%%%%%%%%%%%%%%%%%%%%%%%%%%%%%%%%%%%%%%%

\begin{frame}
 \frametitle{Bug reporting empirical studies}
 \begin{itemize}
    \item Reporters responses about items difficulty
        \begin{itemize}
        \item Observed behavior
        \item Expected behavior
        \item How to reproduce the bug (steps)
        \end{itemize}
 \end{itemize}
\end{frame}

%%%%%%%%%%%%%%%%%%%%%%%%%%%%%%%%%%%%%%%%%%%%%%%%%%%%%%%%%%%%%%

\begin{frame}
 \frametitle{Bug reporting empirical studies}
 \begin{itemize}
    \item Mismatching between developers and reporters?
 \end{itemize}
\end{frame}

%%%%%%%%%%%%%%%%%%%%%%%%%%%%%%%%%%%%%%%%%%%%%%%%%%%%%%%%%%%%%%

\begin{frame}
 \frametitle{Bug reporting empirical studies}
 \begin{itemize}
    \item Developers and reporters agreed in top most important items
    \item Also agreed in less important item
    \item Many mismatching in the rest of the items
    \item Stack traces, test cases and code examples are rarely reported but considered important by developers
 \end{itemize}
\end{frame}

%%%%%%%%%%%%%%%%%%%%%%%%%%%%%%%%%%%%%%%%%%%%%%%%%%%%%%%%%%%%%%

\begin{frame}
 \frametitle{Bug Reporting Info Mistmatching }
 \begin{itemize}
    \item Information used by Developers 
    \item Information provided by Reporters
    \item Spearman Correlation factor: 0.321 (Not really matching!!)
 \end{itemize}
\end{frame}

%%%%%%%%%%%%%%%%%%%%%%%%%%%%%%%%%%%%%%%%%%%%%%%%%%%%%%%%%%%%%%

\begin{frame}
\frametitle{Bug Reporting Mismatching: Helpful for Developers VS Often provided by Reporters }
 \begin{itemize}
    \item Information helpful for Developers 
    \item Information often provided by Reporters
    \item Spearman Correlation factor: -0.035
 \end{itemize}
\end{frame}

%%%%%%%%%%%%%%%%%%%%%%%%%%%%%%%%%%%%%%%%%%%%%%%%%%%%%%%%%%%%%

\begin{frame}
 \frametitle{Bug Reporting Info Mistmatching }
 \begin{itemize}
    \item Information helpful for Developers 
    \item Information expected to be helpful by Reporters
    \item Spearman Correlation factor: 0.839
 \end{itemize}
\end{frame}

%%%%%%%%%%%%%%%%%%%%%%%%%%%%%%%%%%%%%%%%%%%%%%%%%%%%%%%%%%%%%%


\begin{frame}
 \frametitle{Bug reports fixed sooner}
 \begin{itemize}
    \item Bug reports with stack traces
    \item Bug reports easier to read
    \item Bug reports including code samples
 \end{itemize}
\end{frame}

%%%%%%%%%%%%%%%%%%%%%%%%%%%%%%%%%%%%%%%%%%%%%%%%%%%%%%%%%%%%%%

\begin{frame}
 \frametitle{Bug reporting empirical studies}
 \begin{itemize}
    \item Duplicate bug reports
       \begin{itemize}
       \item Demmand more effort from QA teams?
       \item Do developers suffer them?
       \item Do they provide more/useful information?
       \end{itemize}
 \end{itemize}
\end{frame}

%%%%%%%%%%%%%%%%%%%%%%%%%%%%%%%%%%%%%%%%%%%%%%%%%%%%%%%%%%%%%%
\begin{frame}
 \frametitle{Bug reporting empirical studies}
 \begin{itemize}
    \item Duplicate bug reports
       \begin{itemize}
       \item A common practice is to simply close the bug and discard the info.
       \item Is the first one the best one?
       \item Duplicate bugs are not really a problem
       \item Often add more information
       \item Master report (Original one)
       \item Extended report (Master + duplicates)
       \item Extended reports can provide much more useful information
       \end{itemize}
 \end{itemize}
\end{frame}

%%%%%%%%%%%%%%%%%%%%%%%%%%%%%%%%%%%%%%%%%%%%%%%%%%%%%%%%%%%%%%

\section{Mozilla QA and testing tools}

%%%%%%%%%%%%%%%%%%%%%%%%%%%%%%%%%%%%%%%%%%%%%%%%%%%%%%%%%%%%%%

\begin{frame}
 \frametitle{Mozilla QA Tools: Bugzilla}
 \begin{itemize}
 \item Web based bug tracker originally developed to be used by the Mozilla project
 \item Released under the MPL
 \item Released as FLOSS in 1998 by Netscape
 \item It is used by many important FLOSS projects/companies
    \begin{itemize}
    \item GhostScript
    \item Novell
    \item Songbird
    \item Gnome
    \item Apache
    \item ...
    \end{itemize}
 \end{itemize}
\end{frame}

%%%%%%%%%%%%%%%%%%%%%%%%%%%%%%%%%%%%%%%%%%%%%%%%%%%%%%%%%%%%%%
\begin{frame}
\frametitle{Bugzilla Usage}
\begin{center}
 \includegraphics[height=6cm]{figs/Bugzilla_usage_worldwide.png}
\begin{figure}
\end{figure}
\end{center}
\end{frame}

%%%%%%%%%%%%%%%%%%%%%%%%%%%%%%%%%%%%%%%%%%%%%%%%%%%%%%%%%%%%%%

\begin{frame}
\frametitle{Bugzilla Install list web page}
\begin{center}
\begin{figure}
 \includegraphics[height=6cm]{figs/Bugzilla_Install_List.png}
\end{figure}
\end{center}
\end{frame}

%%%%%%%%%%%%%%%%%%%%%%%%%%%%%%%%%%%%%%%%%%%%%%%%%%%%%%%%%%%%%%

\begin{frame}
\frametitle{Mozilla Buzgzilla Web}
\begin{center}
\begin{figure}
 \includegraphics[height=6cm]{figs/MozillaBugzillaWeb.png}
\end{figure}
\end{center}
\end{frame}

%%%%%%%%%%%%%%%%%%%%%%%%%%%%%%%%%%%%%%%%%%%%%%%%%%%%%%%%%%%%%%

\begin{frame}
 \frametitle{Mozilla QA Tools: Bugzilla. How to report a bug}
 \begin{itemize}
    \item Point your browser to \textit{http://bugzilla.mozilla.org}
    \item Recomended but not mandatory, to have a bugzilla acccount.
    \item Log into bugzilla.
    \item Use the search box in the main page to see if your bug has already been submitted.
    \begin{itemize}
       \item Important to avoid duplicate bugs.
    \end{itemize}
    \item If a previous bug exist, you should add your info into as a comment.
    \item If you are sure that your issue has not already been reported then click in the link New in the toolbar on the top or click the File a bug image.
    \item Select the product in wich you want to fill the bug
 \end{itemize}
\end{frame}

%%%%%%%%%%%%%%%%%%%%%%%%%%%%%%%%%%%%%%%%%%%%%%%%%%%%%%%%%%%%%%

\begin{frame}
 \frametitle{Mozilla QA Tools: Bugzilla. How to report a bug}
 \begin{itemize}
    \item Information needed to report a bug (1/2):
    \begin{itemize}
      \item Component. (if you don't know, select 'General').
      \item Version. Of the software that you are using or about you will fill a bug.
      \item Hardware. The architecture of the machine in wich you are running the software.
      \item OS (running the software) Bugzilla will try to guess it!
      \end{itemize}
 \end{itemize}
\end{frame}


\begin{frame}
 \frametitle{Mozilla QA Tools: Bugzilla. How to report a bug}
 \begin{itemize}
    \item Information needed to report a bug (2/2):
    \begin{itemize}
      \item Summary. A short one-sentence summary of the issue.
      \item Description. Full details of the issue, giving as much detail as possible.
      \item Add and Attachment. Use this button if you have important docs to attach to the bugs, like images or something.
      \item Remember to add all the info that you can, the rest can be left blank).
      \item Once you are done click 'Submit Bug'.
      \end{itemize}
 \end{itemize}
\end{frame}

%%%%%%%%%%%%%%%%%%%%%%%%%%%%%%%%%%%%%%%%%%%%%%%%%%%%%%%%%%%%%%

\begin{frame}
 \frametitle{Mozilla QA Tools: Bugzilla. How to report a bug}
 \begin{itemize}
     \item The simple way to report a bug with the right information should be enough.
     \item Interesting essay about how to fill a bug at 
         \begin{itemize}
         \item \textit{http://www.chiark.greenend.org.uk/~sgtatham/bugs.html}
         \end{itemize}
 \end{itemize}
\end{frame}

%%%%%%%%%%%%%%%%%%%%%%%%%%%%%%%%%%%%%%%%%%%%%%%%%%%%%%%%%%%%%%
\begin{frame}
 \frametitle{Mozilla QA Tools: Crash Reporter}
 \begin{itemize}
 \item Tools that sends information to Thunderbird developers whenever Thunderbird crashes
 \item Is a set of libraries for client-side crash reporting
 \item Same thing, different names
    \begin{itemize}
     \item QFA
     \item TalkBack
     \item BreakPad
     \item Socorro
     \item ...
    \end{itemize}
 \end{itemize}
\end{frame}

%%%%%%%%%%%%%%%%%%%%%%%%%%%%%%%%%%%%%%%%%%%%%%%%%%%%%%%%%%%%%%

\begin{frame}
 \frametitle{Crash Reporter Implementation}
 \begin{itemize}
   \item Client integration to handle crash reporting 
   \item Server (Socorro) that collects and processes collected crash data
   \item A web interface (Socorro) for viewing and parsing crash reports, called
   \item This info is used by Thunderbird developers to make future versions crash less frequently. 
   \item Important part of the Mozilla Core
   \item It cannot be deselected during custom install, unlike the old version, the Talkback add-on.
   \item Also used in other Mozilla products (Firefox, etc).
   \item Development started by Google.
 \end{itemize}
\end{frame}

%%%%%%%%%%%%%%%%%%%%%%%%%%%%%%%%%%%%%%%%%%%%%%%%%%%%%%%%%%%%%%

\begin{frame}
 \frametitle{How to use the crash reporter}
 \begin{itemize}
  \item The Crash Reporter appears only when Thunderbird crashes.
  \item Only when Thunderbird crashes, it doesn't cause Thunderbird to crash!, like some people think.
  \item Once it has appeared, it saves a binary "dump" file, with a lot of info about the crash.
  \item Fill the info in the form to tell Mozilla about this crash so they can fix it.
  \item Mozilla Crash Reporter will send a summary of the crash to Mozilla.
  \item Restart your application!
  \item Referencing crash reports in bugzilla is a good idea, just add the GGUID to the bug, when reporting a bug that cause a crash!
 \end{itemize}
\end{frame}

%%%%%%%%%%%%%%%%%%%%%%%%%%%%%%%%%%%%%%%%%%%%%%%%%%%%%%%%%%%%%%

\begin{frame}
\frametitle{Crash Reporter}
\begin{center}
 \includegraphics[height=6cm]{figs/MozillaCrashReporter.png}
\begin{figure}
\end{figure}
\end{center}
\end{frame}

%%%%%%%%%%%%%%%%%%%%%%%%%%%%%%%%%%%%%%%%%%%%%%%%%%%%%%%%%%%%%%

\begin{frame}
\frametitle{Crash Reporter}
\begin{center}
 \includegraphics[height=6cm]{figs/CrashReportContents.png}
\begin{figure}
\end{figure}
\end{center}
\end{frame}

%%%%%%%%%%%%%%%%%%%%%%%%%%%%%%%%%%%%%%%%%%%%%%%%%%%%%%%%%%%%%%

\begin{frame}
\frametitle{Crash reports: Web view}
 \begin{itemize}
  \item Using Web Interface
     \begin{itemize}
       \item Crash Reporter will save info about all the crashes that happend in your software.
       \item In case of Thunderbird we need a extra extension, view:about (Need to install)
       \item about:crashes will show a list with links to all the crashes that your applications has suffered.
       \item Click in one of the links to view the report.
     \end{itemize}
 \end{itemize}
\end{frame}

%%%%%%%%%%%%%%%%%%%%%%%%%%%%%%%%%%%%%%%%%%%%%%%%%%%%%%%%%%%%%%

\begin{frame}
\frametitle{Crash Reporter Web View}
\begin{center}
 \includegraphics[height=6cm]{figs/CrashWebView.png}
\begin{figure}
\end{figure}
\end{center}
\end{frame}

%%%%%%%%%%%%%%%%%%%%%%%%%%%%%%%%%%%%%%%%%%%%%%%%%%%%%%%%%%%%%%

\begin{frame}
\frametitle{Crash reports: File view}
 \begin{itemize}
 \item Saved a local copy on your system
  \item Locations (OS Depending)
    \begin{itemize}
      \item Linux: .Mozilla/firefox/Crash Reports/
    \end{itemize}
  \item Each submitted crash report is identified as a text file, located in the Crash Reports folder (within the submitted subfolder).
  \item Find the GGUID of the crash report that we want to take a look at...
  \item ...append the GGUID to the end of \textit{http://crash-stats.mozilla.com/report/index/}
 \end{itemize}
\end{frame}

%%%%%%%%%%%%%%


\begin{frame}
\frametitle{Crash Reporter: File View}
\begin{center}
 \includegraphics[height=6cm]{figs/MozillaCrashReportFile.png}
\begin{figure}
\end{figure}
\end{center}
\end{frame}

%%%%%%%%%%%%%%%%%%%%%%%%%%%%%%%%%%%%%%%%%%%%%%%%%%%%%%%%%%%%%%

\begin{frame}
 \frametitle{Bug Reporting Hands on!}
 \begin{itemize}
    \item Let's try to crash it!! 
    \item Create an account in Mozilla Bugzilla system (need e-mail confirmation)
        \begin{itemize}
        \item https://bugzilla.mozilla.org/
        \end{itemize}
    \item Try to crash Thunderbird or Firefox reproducing an existing bug
        \begin{itemize}
        \item You can look for bug reports in Bugzilla database
        \end{itemize}
    \item Send the crash report
    \item View the crash report information (Web and File View)
 \end{itemize}
\end{frame}

%%%%%%%%%%%%%%%%%%%%%%%%%%%%%%%%%%%%%%%%%%%%%%%%%%%%%%%%%%%%%


\begin{frame}
\frametitle{Litmus}
 \begin{itemize}
    \item Litmus is a web-based, open source QA tool maintained by Mozilla Corporation.
    \item Is designed to improve workflow, visibility, and turnaround time in the Mozilla QA process.
    \item Everything is manual.
 \end{itemize}

\end{frame}

%%%%%%%%%%%%%%%%%%%%%%%%%%%%%%%%%%%%%%%%%%%%%%%%%%%%%%%%%%%%%%
\begin{frame}
\frametitle{Litmus}
 \begin{itemize}
    \item Litmus does:
       \begin{itemize}
       \item Makes easier for casual testers to assist with testing Mozilla products.
       \item Serve as a repository for test cases, with all the inherent management abilities that implies.
       \item Serve as a repository for test results, carrying over the best features of Testrunner, e.g. test lists, division of labor, etc.
       \item Provide a query interface for viewing, reporting on, and comparing test results.
       \item Show a web services interface for the mechanical batch submission of testing results.
       \end{itemize}
 \end{itemize}
\end{frame}

%%%%%%%%%%%%%%%%%%%%%%%%%%%%%%%%%%%%%%%%%%%%%%%%%%%%%%%%%%%%%%

\begin{frame}
\frametitle{Litmus}
 \begin{itemize}
    \item Litmus does not:
       \begin{itemize}
       \item Manage the automation of testing.
       \item At all! :)
       \item There are some other tools for that in Mozilla:
          \begin{itemize} 
          \item Mozmill
          \end{itemize}
       \end{itemize}
 \end{itemize}

\end{frame}


%%%%%%%%%%%%%%%%%%%%%%%%%%%%%%%%%%%%%%%%%%%%%%%%%%%%%%%%%%%%%%

%%%%%%%%%%%%%%%%%%%%%%%%%%%%%%%%%%%%%%%%%%%%%%%%%%%%%%%%%%%%%%

\begin{frame}
\frametitle{Test types}
  \begin{itemize}

  \item Smoketest 
       \begin{itemize}
       \item Covers common operations (Install, Uninstall, Usability)
       \item Includes a set of various operations
       \item Its the Standard Test.
       \end{itemize}
  \item Full Functional Tests
       \begin{itemize}
       \item Its the biggest testrun in Litmus. 
       \item It includes a lot of testcases for all aspects of the Product you have selected.
       \end{itemize}
  \item Localization (l10)
       \begin{itemize}
       \item A test provided for Localization Builds
       \end{itemize}
  \item Basic Functional Test (BFT)
       \begin{itemize}
       \item Includes more testcases than a Smoketest and covers all aspects of the Product.
       \item From Installing till Security Features.
       \end{itemize}

  \end{itemize}
\end{frame}

%%%%%%%%%%%%%%%%%%%%%%%%%%%%%%%%%%%%%%%%%%%%%%%%%%%%%%%%%%%%%%

\begin{frame}
\frametitle{Test results}
 \begin{itemize}
   \item Not Run: The test has not been excuted yet
   \item Pass: Test got the expected results
   \item Fail: Test failed
   \item Test unclear/broken: Not enought information to run the test
 \end{itemize}
\end{frame}

%%%%%%%%%%%%%%%%%%%%%%%%%%%%%%%%%%%%%%%%%%%%%%%%%%%%%%%%%%%%%%

\begin{frame}
\frametitle{Testing...}
 \begin{itemize}
   \item We need a user account in Litmus
   \item A testing build of Thunderfird
   \item Set up for the test run
   \item Execute the test
   \item Submit the test run results
 \end{itemize}
\end{frame}

%%%%%%%%%%%%%%%%%%%%%%%%%%%%%%%%%%%%%%%%%%%%%%%%%%%%%%%%%%%%%%
\begin{frame}
\frametitle{Manual Testing with Litmus}
\begin{center}
 \includegraphics[height=6cm]{figs/Litmus_01_Config_Required_Criteria.png}
\begin{figure}
\end{figure}
\end{center}
\end{frame}

%%%%%%%%%%%%%%%%%%%%%%%%%%%%%%%%%%%%%%%%%%%%%%%%%%%%%%%%%%%%%%

\begin{frame}
\frametitle{Manual Testing with Litmus}
\begin{center}
 \includegraphics[height=6cm]{figs/Litmus_02_Testgroup_Selection.png}
\begin{figure}
\end{figure}
\end{center}
\end{frame}

%%%%%%%%%%%%%%%%%%%%%%%%%%%%%%%%%%%%%%%%%%%%%%%%%%%%%%%%%%%%%%

\begin{frame}
\frametitle{Manual Testing with Litmus}
\begin{center}
 \includegraphics[height=6cm]{figs/Litmus_03_Test_Run_Intro.png}
\begin{figure}
\end{figure}
\end{center}
\end{frame}

%%%%%%%%%%%%%%%%%%%%%%%%%%%%%%%%%%%%%%%%%%%%%%%%%%%%%%%%%%%%%%

\begin{frame}
\frametitle{Manual Testing with Litmus}
\begin{center}
 \includegraphics[height=6cm]{figs/Litmus_04_Test_Results_Submitted.png}
\begin{figure}
\end{figure}
\end{center}
\end{frame}

%%%%%%%%%%%%%%%%%%%%%%%%%%%%%%%%%%%%%%%%%%%%%%%%%%%%%%%%%%%%%%

\begin{frame}
 \frametitle{Manual Testing Hands on!}
 \begin{itemize}
    \item Let's test Thunderbird!! 
    \item Create an account in Litmus
        \begin{itemize}
        \item https://litmus.mozilla.org/
        \end{itemize}
    \item Perform a simple test run 
        \begin{itemize}
        \item Follow the instructions
        \item Rate the test
        \end{itemize}
    \item Submit the test results
 \end{itemize}
\end{frame}

%%%%%%%%%%%%%%%%%%%%%%%%%%%%%%%%%%%%%%%%%%%%%%%%%%%%%%%%%%%%%%


\begin{frame}
\frametitle{Mozmill}
 \begin{itemize}
  \item Test tool and framework for writing automated tests for Gecko based applications 
     \begin{itemize}
     \item Firefox
     \item Thunderbird
     \item Songbird, etc.
     \end{itemize}
  \item Build as an extension (including IDE) or as a command line client
  \item Mozmill is a test tool to run, debug and write tests
 \end{itemize}
\end{frame}

%%%%%%%%%%%%%%%%%%%%%%%%%%%%%%%%%%%%%%%%%%%%%%%%%%%%%%%%%%%%%%

\begin{frame}
\frametitle{Mozmill}
 \begin{itemize}
  \item Used by the Firefox QA and Thunderbird QA teams to automate test
  \item Easy to install and collaborate
  \item Tests are written in JavaScript
  \item Have access to any part of the UI and also to all available XPCOM components
      \begin{itemize}
      \item Cross Platform Component Object Model
      \item Similar to CORBA
      \end{itemize}
 \end{itemize}
\end{frame}

%%%%%%%%%%%%%%%%%%%%%%%%%%%%%%%%%%%%%%%%%%%%%%%%%%%%%%%%%%%%

\begin{frame}
\frametitle{Mozmill Test Automation}
\begin{center}
 \includegraphics[height=6cm]{figs/Mozmill_Test_Automation_Map.png}
\begin{figure}
\end{figure}
\end{center}
\end{frame}

%%%%%%%%%%%%%%%%%%%%%%%%%%%%%%%%%%%%%%%%%%%%%%%%%%%%%%%%%%%%%%
%%% http://piratepad.net/zyyP4uNTsp
%%%%%%%%%%%%%%%%%%%%%%%%%%%%%%%%%%%%%%%%%%%%%%%%%%%%%%%%%%%%%%

\section{References and Links}

\begin{frame}
 \frametitle{There are many other FLOSS QA Tools}
 \begin{itemize}
    \item Selenium
        \begin{itemize}
        \item Simple test automation for web applications
        \item Easy and quick
        \end{itemize}
    \item WebDriver
        \begin{itemize}
        \item Google testing framework
        \item Test can be written in many languages
        \end{itemize}
    \item GNU/LDTP
        \begin{itemize}
        \item Linux Desktop Testing Project 
        \item Accessibility Libraries
        \end{itemize}
  \end{itemize}

\end{frame}

%%%%%%%%%%%%%%%%%%%%%%%%%%%%%%%%%%%%%%%%%%%%%%%%%%%%%%%%%%%%%%

\begin{frame}
 \frametitle{References}
 \begin{itemize}
    \item SuMoMo \textit{http://support.mozillamessaging.com}
    \item SuMo \textit{http://support.mozilla.com}
    \item bugzilla \textit{http://bugzilla.libresoft.es}
    \item breakpad \textit{http://code.google.com/p/google-breakpad/}
    \item crash-stats \textit{http://crash-stats.mozilla.com/report/index/}
    \item Wiki Mozilla \textit{http://wiki.mozilla.org}
    \item OpenQA \textit{http://openqa.org}
  \end{itemize}

\end{frame}

%%%%%%%%%%%%%%%%%%%%%%%%%%%%%%%%%%%%%%%%%%%%%%%%%%%%%%%%%%%%%%


\end{document}
